\section{Week 3 Generated \texorpdfstring{$\mu$}{mu} Control}
\label{sec:week3-generated-mu}

\paragraph{Research question and claim.}
Week~3 targets \claimid{OBL-MU-TOPLEVEL-CONTROL}. The required statement is no
longer about proof-only local wrappers; the actual generated procedures must
appear directly in the theorem:
\[
  \take_{32}\!\bigl(
    \mathsf{sf\_mu\_rawpre}(sk,pre,m)
  \bigr)
  =
  \mathsf{verify\_hash\_mu}(vk,pre,m).
\]
More precisely, the claim is
\[
  \mathsf{raw\_prelen}(\ell)\land
  \prefix_{992}(sk,mem,vk)
  \Longrightarrow
  \take_{32}(
    \mathsf{sf\_mu\_rawpre}(sk,pre,m)
  )
  =
  \mathsf{verify\_hash\_mu}(vk,pre,m).
\]

\paragraph{Generated control sequence.}
The Sign side is the actual generated
\procname{SignMuHashTarget.M.\_sf\_mu\_rawpre}. Its control sequence is:
\begin{enumerate}
  \item stack/local initialization and zeroed Keccak state;
  \item \procname{\_keccak\_init\_state};
  \item \procname{\_sf\_shake256\_absorb\_sk} with \texttt{vkbytes = 992};
  \item raw absorb of \texttt{pre};
  \item raw absorb of \texttt{message};
  \item \procname{\_sf\_shake256\_finalize};
  \item one additional \procname{\_keccakf1600};
  \item \procname{protect\_ptr};
  \item \procname{\_sf\_shake256\_squeeze64}.
\end{enumerate}
The Verify side is the actual generated
\procname{VerifyMuHashTarget.M.\_\_verify\_hash\_mu}. After initialization and
the 992-byte VK absorb, it computes the branch guard
\[
  ctxflag \eqdef prelen \gg 63.
\]
The raw branch is taken only after this actual generated shift computes zero.

\paragraph{Why exact-to-wrapper refinement was chosen.}
The Week~2 helper theorem already proved the semantic core of the raw hash
path, but only through local wrappers that reconstructed the top-level control
shape. A direct \texttt{inline; sim} proof over the two generated procedures
stalled on:
\begin{itemize}
  \item generated local-variable scheduling;
  \item erased spill/unspill/declassify instructions;
  \item the Verify raw/context branch split; and
  \item mismatch between the generated tail-call boundary and the wrapper-level
        theorem boundary.
\end{itemize}
Week~3 therefore isolates two unary exact adapters and composes them by
transitivity rather than attempting one monolithic whole-program simulation.

\paragraph{Compiled exact adapters.}
Theory \theoryname{ExactMuTopControl} proves:
\begin{itemize}
  \item \theoryname{sf\_mu\_rawpre\_refines\_raw\_top};
  \item \theoryname{raw\_prelen\_shr63\_zero};
  \item \theoryname{raw\_prelen\_branch\_guard}; and
  \item \theoryname{verify\_hash\_mu\_raw\_refines\_raw\_top}.
\end{itemize}
The Verify-side predicate is
\[
  \mathsf{raw\_prelen}(prelen)
  \eqdef
  0 \le \mathsf{to\_uint}(prelen) < 2^{63}.
\]
The two branch lemmas prove that this predicate implies that the actual
generated shift-by-63 word operation yields \(0\), so the raw branch is
selected semantically rather than by assumption.

\paragraph{Actual top-level theorem.}
The closing theorem is
\theoryname{sign\_verify\_generated\_raw\_mu\_prefix}. Its left and right
procedures are exactly:
\begin{verbatim}
Sign._sf_mu_rawpre
Verify.__verify_hash_mu
\end{verbatim}
Its premises include:
\begin{itemize}
  \item Sign \texttt{vkbytes = 992};
  \item Verify \texttt{vklen = 992};
  \item equal raw pre/message pointers and lengths;
  \item \theoryname{raw\_prelen prelen};
  \item a KeyGen-derived \theoryname{sk\_memory\_prefix}; and
  \item the explicit VK, pre, and message no-wrap predicates, including
        \(\mathsf{to\_uint}(base)+992<2^{64}\).
\end{itemize}
Its conclusion is bytewise equality of the first 32 Sign output bytes and the
entire Verify output.

\paragraph{Reachability and non-vacuity.}
The theorem is not justified by a generic witness over unrelated arrays.
Theory \theoryname{ExactMode2RawMuComposition} proves
\theoryname{generated\_raw\_mu\_preconditions\_from\_keygen\_prefix}, which
packages:
\begin{itemize}
  \item the returned KeyGen prefix theorem;
  \item the concrete witness \(base = 0\);
  \item \theoryname{raw\_prelen W64.zero}; and
  \item the constructed Verify memory image obtained by storing the packed VK.
\end{itemize}
It further proves the Hoare theorem
\theoryname{keypair\_internal\_return\_reaches\_generated\_raw\_mu\_preconditions}
over the actual generated
\procname{crypto\_sign\_keypair\_internal\_mode2\_jazz}. Every terminating
execution therefore supplies this constructed local-array/raw premise set;
KeyGen termination is neither assumed nor concluded. Thus the actual
generated theorem premises are satisfiable by a concrete call path.

\paragraph{Status.}
\claimid{OBL-MU-TOPLEVEL-CONTROL} is \status{PROVED} for the raw/internal path.
The public-context branch remains separate and unproved.
